\section{The Agent Layer as a Design Position}\label{sec:agents}

\subsection{Verdicts computed, prose generated}
Five agents (document extraction, order entry, exception triage, invoice audit and an analytics assistant) share one rule: every verdict is computed, and a language model writes only the sentence a person reads. Severity is arithmetic over how far past its threshold a leg is predicted to run; an invoice verdict is four comparisons. The rule makes each decision reproducible and therefore scorable, and bounds what a model error can do to a sentence. It follows the division of labour of systems in which a language model explains an optimiser's output while the optimiser decides~\cite{li2023optiguide}, applied more strictly, and builds on tool-using agent designs~\cite{yao2023react,wang2024survey,xi2023rise} exposed through a standard protocol~\cite{mcp2025specification}. Each agent's evaluation set was authored by someone other than its builder, and every result is reported beside a trivial policy on the same set. All agent data are synthetic or generated for this study; only the network data are real.

\paragraph{Scale re-runs.} First sets of 20 to 50 cases give perfect scores that are consistent with a true rate of 84\%, so the agents were re-run at scale, stratified so that hard cases form a named share, with Wilson intervals~\cite{wilson1927probable,agresti1998approximate}. Order entry scored \textbf{\fzOrderEvalScaleHeldoutPct{}\% on \fzOrderEvalScaleCases{} held-out emails} (95\% interval [79.4, 89.3]), falling to \fzOrderEvalScaleAdversarialHeldoutPct{}\% on the adversarial half, where every error was an order filed on a value the email never stated. Invoice audit was correct on \textbf{\fzInvoiceEvalScaleAccuracyPct{}\% of \fzInvoiceEvalScaleCases{} invoices} (interval [85.1, 93.4]) with \fzInvoiceEvalScaleFalseDisputes{} correct invoices disputed; its misses are the one kind that falls inside its own tolerance. The orchestrator routed \textbf{\fzLifecycleEvalScaleRoutingPct{}\% of \fzLifecycleEvalScaleCases{} lifecycles} as the email deserved, exercising all \fzLifecycleEvalScalePaths{} paths of its graph. Document extraction scored \fzDocExtractionAccuracyPct{}\% of fields on \fzDocExtractionRows{} of 40 planned rows (the remaining rows were not run for quota reasons). The assistant routed \fzAssistantRouteAccuracyLlmPct{}\% of 30 questions correctly with model phrasing and \fzAssistantRouteAccuracyPct{}\% without.

\paragraph{The 100\% on the 200-email set is not held out.} The rule-based order agent was fixed after the 200-email set found its first failures, so its later 100\% on that set is a fitted score. We therefore built a new held-out hard set (\fzHardOrderCases{} emails over hard templates such as Hinglish quantities, conflicting figures and quoted threads) and scored the \emph{unmodified} agents once: the order agent scored \textbf{\fzHardOrderRulesPct{}\%}, failing three template classes outright; the invoice auditor scored \fzHardInvoiceRulesPct{}\% on a boundary set. We did not repair the order failures, which are documented.

\subsection{Exception precision with uncertainty}\label{sec:exception}
Section~\ref{sec:replay} gives the as-of precision. The comparison with alerting on every leg is the one the replay's first version exaggerated: the corrected gain is \fzExcGainPointsDpTwo{} points, significant under the valid paired test, and the lowest severity grade performs worse than alerting everything.

\subsection{What the language model adds: a pilot ablation}\label{sec:ablation}
The design position makes a direct question answerable: what does each agent gain or lose when a model makes the decision itself? We ran the order agent with an open 120-billion-parameter model deciding, five sampled runs per case plus two at temperature zero, on the same stratified cases the rules saw. \emph{This is a pilot}: the free-tier daily token cap bound before the pre-registered 40 cases per agent, so the sample is \fzAblationOrderCases{} order cases (and 8 each for invoice and exception triage), a second small model at 8 each is reported separately, and the full design is resumable but not run. On those cases the rules scored 100\% (not held out, see above) and the model deciding scored \fzAblationOrderLlmAccuracyPct{}\% pooled over five runs (95\% interval [\fzAblationOrderLlmCiLowPct{}, \fzAblationOrderLlmCiHighPct{}]); the interval of the difference to the rules excludes zero. The model was less reproducible (all five runs agreed on \fzAblationOrderUnanimityPct{}\% of cases; Fleiss $\kappa=\fzAblationOrderKappaDpTwo{}$~\cite{fleiss1971measuring}) and filed 12\% of its orders on values the email did not state, against none for the rules. On invoice audit and severity grading, where the policy is arithmetic and given in the prompt, the pilot found no difference at $n=8$. The honest summary is that the evidence supports ``comparable where the policy is arithmetic'' and ``worse where it needs reading'' for the order task on this sample, and does not support the stronger claim that a deterministic core is generally better or cheaper: cost on a paid tier would be about a quarter of a dollar per thousand verdicts at assumed list prices that we did not verify, and latency is seconds against microseconds. The empty column of the study is human labels: no agent decision has been compared against a second labeller, and the protocol and sheets for that are prepared but not run.
